\chapter{Create a GitHub Account}

Creating a GitHub account is quite simple. However, if you are creating a secondary account, you must be logged out of your primary account in the browser, otherwise, things will get confusing! To sign out:

\begin{itemize}
    \item Click your profile picture in the top right corner.
    \item Click ``Sign Out''
    \item Confirm by clicking one or other of the subsequent ``sign out'' buttons.
\end{itemize}

Ok, the slate is clean, lets create an account.

\begin{itemize}
    \item Point your browser at \url{https://github.com}.
    \item Click ``Sign up''.
    \item On the right:
    \item Fill in your email address.
    \item Fill in an initial password, be secure! You need either 15 characters, or, 8 characters which must include at least one digit and at least one lowercase letter.
    \item Fill in a username. Replace spaces with hyphens---but only one consecutive hyphen is allowed. You will be told if the username is available or not.
    \item Select a country.
    \item Tick the box for email updates if you want them. You don't have to though. Your choice.
    \item Click ``Create account''.
\end{itemize}

You are now prompted for a nine digit code. Check the email address that you supplied and enter the digits. Once you have done so, you will need to login with the username (or email) and password you supplied above.

Job done, you can now create repositories---see \chapref{chap: Creating Repositories}---and the like. Read on, however, for details of how to use \emph{https} and \emph{ssh} protocols to be able to push changes back to a repository.
